% !PRINTED 574
\weberpage{574}
%<<CONT
worin die Kongruenz auf die Perioden bezogen wird, also (\S\ 152, 1.)
\weq{2}{\begin{array}{rll}
\text{im Falle a),} & \text{wenn } \nu \equiv \pm\nu', & \\
\text{b),} & \text{wenn } \nu \equiv \pm\nu', \ \pm i\nu' & (\mathrm{mod}\ \mu), \\
\text{c),} & \text{wenn } \nu \equiv \pm\nu', \ \pm\varrho\nu', \ \pm\varrho^2\nu'. &
\end{array}}
Es seien jetzt
\weq{3}{\mu = \mathfrak{a}\,\mathfrak{m}, \quad \mu_1 = \mathfrak{a}_1\,\mathfrak{m}}
zwei ganze Zahlen mit dem größten gemeinschaftlichen Idealteiler $\mathfrak{m}$ im Körper $\Omega$. Dann haben $T_\mu$ und $T_{\mu_1}$ einen gemeinsamen Linearteiler, wenn
\weq{4}{\tau\left(\frac{\nu}{\mu}\right) = \tau\left(\frac{\nu_1}{\mu_1}\right),}
und dies findet nach (2) dann und nur dann statt, wenn
\weq{5}{\mu\nu_1 \equiv \varepsilon\mu_1\nu \quad (\mathrm{mod}\ \mu\mu_1),}
worin $\varepsilon$ eine \gesperrt{Einheit} des Körpers $\Omega$ ist, also $\varepsilon = \pm 1$ im allgemeinen, $\varepsilon = \pm 1, \pm i$, wenn $\Delta = -4$, und $\varepsilon = \pm 1, \pm\varrho, \pm\varrho^2$, wenn $\Delta = -3$ ist. Aus (5) folgt aber, daß $\nu$ durch $\mathfrak{a}$ und $\nu_1$ durch $\mathfrak{a}_1$ teilbar sein muß.

Denken wir uns $\nu$ gegeben und $\nu_1$ gesucht, so ist die Kongruenz (5) nur möglich, wenn $\mu_1\nu$ durch $\mu$ teilbar ist, und da $\mathfrak{a}$ und $\mathfrak{a}_1$ relativ prim sind, so muß $\nu$ durch $\mathfrak{a}$ teilbar sein. Ist $\alpha$ eine durch $\mathfrak{a}$ teilbare ganze Zahl in $\Omega$, so beschaffen, daß $\alpha : \mathfrak{a}$ relativ prim zu $\mathfrak{m}$ ist, so können wir demnach
\weq{6}{\nu \equiv \alpha\xi \quad (\mathrm{mod}\ \mu)}
setzen (Bd. II, \S\ 166, 7.) und erhalten die sämtlichen voneinander verschiedenen Werte $\tau\left(\dfrac{\nu}{\mu}\right)$, und jeden zwei-, vier- oder sechsmal, wenn wir $\xi$ ein volles Restsystem nach dem Modul $\mathfrak{m}$ durchlaufen lassen (mit Ausschluß der Null). Man erhält dann die sämtlichen gemeinsamen Wurzeln von $T_\mu$ und $T_{\mu_1}$ in der Form
\weq{7}{\tau_\xi = \tau\left(\frac{\alpha\xi}{\mu}\right).}
Der größte gemeinschaftliche Teiler $D_{\mathfrak m}$ von $T_\mu$ und $T_{\mu_1}$ hat also alle die Größen (7) zu Wurzeln. Aus $D_{\mathfrak m}$ läßt sich noch auf rationalem Wege ein Teiler $T_{\mathfrak m}$ absondern, der nur die unter den Größen (7) zu Wurzeln hat, in denen $\xi$ relativ prim zu $\mathfrak{m}$ ist, wenn man $D_{\mathfrak m}$ von allen Faktoren befreit, die es mit einem $D_{\mathfrak m'}$ gemein hat, wenn $\mathfrak{m}'$ ein Teiler von $\mathfrak{m}$ ist. Da in dieser
%CONT>>
